\documentclass[10pt]{article}

\IfFileExists{neurips_2026.sty}{
  \usepackage[preprint]{neurips_2026}
}{
  \usepackage[margin=1in]{geometry}
}
\usepackage{amsmath,amssymb,booktabs,graphicx,microtype,multirow}
\usepackage[numbers,sort&compress]{natbib}
\usepackage{xcolor}
\usepackage{hyperref}
\usepackage{url}

\newcommand{\system}{\textsc{AgentRelay}}
\newcommand{\relaytax}{\mathrm{RelayTax}}
\newcommand{\tbd}{\textcolor{red}{TBD}}
\newcommand{\methodtag}[1]{\textsc{#1}}

\title{\system: Switch-Cost-Aware and Transaction-Safe State Handoffs\\
for Edge--Cloud LLM Agents}
\author{Anonymous Authors}
\date{}

\begin{document}
\maketitle

\begin{center}
\fbox{\parbox{0.94\linewidth}{
\textbf{Internal pre-experimental draft.}
This manuscript specifies falsifiable claims, public-data protocols, and
pre-registered analyses. Every empirical cell is intentionally marked
\tbd{} until produced by native inference under the locked configuration.
The draft must not be submitted or used to claim empirical superiority in
its present form.}}
\end{center}

\begin{abstract}
Edge--cloud language-model agents can reduce monetary cost and latency by
executing routine steps locally and escalating difficult steps to a stronger
cloud model. Existing routers largely treat each decision as model selection:
they estimate which executor is likely to succeed, then choose one under a
cost or quality objective. Stateful agents make this abstraction incomplete.
A switch also transfers a partially executed trajectory, and retries may
repeat external effects. The resulting relay tax---serialization, network
transfer, context reprocessing, validation, and repair---can reverse an
otherwise sensible routing decision. We present \system, a handoff protocol
and joint policy that chooses an executor, a state-transfer mode, and an
effect-commit mode together. \system{} uses typed, content-addressed state
deltas; validates identity, provenance, invariants, and effect histories at
the receiver; requests narrow patches when validation fails; and records
idempotent effect keys in a transaction ledger. A switch-cost-aware selector
uses measured rather than nominal costs and can abstain from switching when
the expected gain is smaller than the relay tax. We define an evaluation on
public ALFWorld, WebShop, AppWorld, and AgentProcessBench data, with immutable
splits, native model inference, recorded network traces, and paired
uncertainty tests. Results are pending. The contribution of this draft is
therefore a complete, auditable hypothesis and implementation package rather
than an unsupported empirical claim.
\end{abstract}

\section{Introduction}
\label{sec:introduction}

Tool-using language-model agents repeatedly observe an environment, reason
about the next step, call tools, and update an evolving trajectory
\citep{yao2023react,schick2023toolformer,liu2024agentbench}. Running every
step on a large remote model is often unnecessary, whereas running every step
on a smaller local model can sacrifice task success. Recent routing systems
select among models or execution sites using predicted quality, confidence,
cost, or latency
\citep{chen2024frugalgpt,ong2025routellm,paul2026agentrouter,
zhang2026hera,son2026swerouter,liu2026agenticrouting}. These systems motivate
step-level edge--cloud execution, but they also expose a systems question that
is easy to hide in a model-selection formulation: what exactly crosses the
boundary when an agent switches executors?

The answer is more than a prompt. At step $t$, the receiving executor may
need the task identity, normalized observations, compact evidence, the current
plan, tool outputs, open obligations, and the status of external effects. A
full replay can be faithful but expensive. A summary can be cheap but omit a
binding constraint. A retry after a timeout may repeat a purchase, message,
or file operation even if model generation itself is deterministic. Thus a
handoff is a state-transfer and transaction event, not merely a new inference
request.

We call the total incremental cost of that event the \emph{relay tax}. It
includes state encoding, bytes transmitted under the observed link, receiver
validation, context prefill, and any repair request. The tax depends on the
current trajectory and transfer mode; it is not a constant surcharge. A
router that ignores it can choose the cloud because the cloud has a higher
predicted success probability, even though staying local has lower measured
cost for the same constrained objective. Conversely, transferring a small
typed delta can make a beneficial switch affordable when full replay cannot.

\system{} makes handoff a first-class protocol. Each decision selects the
compound action
\[
 a_t=(e_t,m_t,c_t),
\]
where $e_t$ is the next executor, $m_t$ is a transfer mode, and $c_t$ is an
effect-commit mode. The state packet is typed and content-addressed. The
receiver verifies task identity, parent state, required fields, provenance,
and legal effect transitions before execution. A failed check triggers a
narrow patch request rather than silently proceeding with incomplete state.
For effectful tools, an idempotency ledger separates authorization,
preparation, commitment, and compensation so that a retry cannot casually
duplicate an already committed action.

This paper advances four testable contributions:
\begin{enumerate}
  \item \textbf{Measured switch-cost routing.} We formulate the relay tax and
  show analytically when omitting it reverses the action selected by a
  constrained router. The implementation records serialization, transfer,
  validation, repair, and context-processing costs per handoff.
  \item \textbf{Joint executor--payload decisions.} We route over executor,
  state-transfer mode, and commit mode rather than selecting a model first and
  treating state movement as fixed plumbing.
  \item \textbf{Typed delta transfer with verified repair.} We introduce
  content-addressed deltas, receiver-side invariants, and patch-on-demand to
  test whether smaller payloads preserve replay-equivalent decision state.
  \item \textbf{Transaction-safe agent effects.} We integrate an effect
  ledger with the handoff protocol and test exactly-once behavior under
  retries, timeouts, and executor switches.
\end{enumerate}

The claims are deliberately narrower than broad agent superiority. We ask
whether measuring handoff cost changes routing, whether joint routing improves
the quality--cost--latency frontier, whether verified deltas preserve state,
and whether the ledger prevents duplicated effects. These questions are
evaluated independently so that a negative answer to one does not disappear
inside an aggregate task-success score.

\section{Background and Problem Formulation}
\label{sec:problem}

\subsection{Stateful edge--cloud agents}

An episode has task $x$, environment observations $o_{1:T}$, tool calls
$u_{1:T}$, and latent decision state $s_t$. An executor belongs to
$\mathcal{E}=\{\mathrm{edge},\mathrm{cloud}\}$, although the formulation
extends to more models. Let $h_t$ denote the materialized trajectory visible
at step $t$. A conventional router estimates an outcome utility $I(e_t\mid
h_t)$ and an inference cost $C_{\mathrm{infer}}(e_t,h_t)$. This misses the
cost and correctness of making $h_t$ usable at a different executor.

We factor the transferable state into five typed components:
\[
s_t=(q_t,o_t,z_t,\pi_t,\ell_t).
\]
Here $q_t$ identifies the immutable task, $o_t$ is normalized world state,
$z_t$ contains evidence and tool outputs with provenance, $\pi_t$ is the
current plan and obligations, and $\ell_t$ is the effect ledger. The sender
may reuse resident state, send a full packet, or send a delta against a
content hash known by the receiver.

\subsection{Measured relay tax}

For candidate action $a=(e,m,c)$ at step $t$, we measure
\begin{align}
H_t(a) ={}&
 \lambda_b B_t(m)
 +\lambda_n L_t^{\mathrm{net}}(B_t(m))
 +\lambda_s L_t^{\mathrm{ser}}(m) \nonumber\\
&+\lambda_v L_t^{\mathrm{verify}}(m)
 +\lambda_p L_t^{\mathrm{prefill}}(e,m)
 +\mathbb{E}\!\left[\lambda_r L_t^{\mathrm{repair}}(m)\right],
\label{eq:relay-tax}
\end{align}
where $B_t$ is transmitted bytes and each $L$ is an observed duration or
resource quantity. The weights put heterogeneous measurements into the
deployment objective. No term is inferred from a nominal bandwidth alone:
the evaluation replays recorded link traces and retains raw component
measurements.

Let $p_t(a)$ be predicted probability of eventual task success, $I_t(a)$ the
direct inference cost, and $R_t(a)$ the expected risk cost of invalid state or
external-effect failure. The unconstrained scalar form is
\begin{equation}
J_t(a)=I_t(a)+H_t(a)+R_t(a)-\beta p_t(a).
\label{eq:objective}
\end{equation}
The deployed selector instead minimizes measured cost subject to
$p_t(a)\geq\tau_q$, a latency budget, memory feasibility, and protocol
validity. Equation~\ref{eq:objective} is useful for analysis and for
displaying Pareto trade-offs; it is not permission to tune $\beta$ on test
episodes.

\paragraph{Switch-cost reversal.}
Consider staying with action $s$ and switching with action $q$. Suppose the
switch has relay tax $H_q>0$ and staying has no handoff tax. A router that
ignores handoff chooses $q$ when
$\beta(p_q-p_s)>I_q-I_s$. The measured objective chooses $s$ when
$\beta(p_q-p_s)<I_q-I_s+H_q$. Therefore the decisions disagree exactly in
the interval
\begin{equation}
 I_q-I_s < \beta(p_q-p_s) < I_q-I_s+H_q.
\label{eq:reversal}
\end{equation}
This elementary interval is operationally important: a larger predicted
model gain need not justify a stateful switch. Experiment E1 measures how
often real trajectories fall in this interval and whether those reversals
improve the realized objective.

\section{Related Work}
\label{sec:related}

\paragraph{LLM and agent routing.}
Cascade and routing work demonstrates that heterogeneous models can improve
the quality--cost frontier. FrugalGPT composes model calls under a budget
\citep{chen2024frugalgpt}; RouteLLM learns preference-based routers
\citep{ong2025routellm}; and AgentRouter specializes routing for agentic tasks
\citep{paul2026agentrouter}. More recent systems move routing inside an
episode. Hera uses hierarchical orchestration for heterogeneous agents
\citep{zhang2026hera}; SWE-Router routes software-engineering steps
\citep{son2026swerouter}; and harness-native Agentic Routing makes the
execution harness aware of model choice \citep{liu2026agenticrouting}.
Speculative and co-serving systems additionally optimize latency and resource
use \citep{ye2026speculative,xiang2026scheduler,dekoninck2025unified}.
Confidence-token routing is another learned uncertainty signal for this
decision \citep{chuang2025confidence}.
\system{} does not claim to originate step-level routing. Its distinction is
to expose state transfer, receiver validation, and effect commitment as
jointly selected and measured parts of each routing action.

\paragraph{Context and state compression.}
Prompt-compression methods reduce token count while attempting to retain
task-relevant information \citep{jiang2023llmlingua,jiang2024longllmlingua}.
Agent-specific work studies folding or compressing long execution state
\citep{sun2026contextfolding,sharma2026statecompression}. These methods can be
used inside a payload encoder, but a compressed string does not by itself
establish that the receiver has the correct task, parent version, evidence
provenance, or effect history. \system{} instead treats a delta as a typed
state transition with explicit invariants. This choice is intentionally less
general than free-form summarization and more directly auditable.

\paragraph{Portable and distributed agents.}
Prior work studies portable context and agent handoff
\citep{ravindran2026portable,kc2026handoff}, placement across edge and cloud
\citep{wang2025placement}, and atomic execution for multi-agent systems
\citep{mohammadi2026atomix}. These lines establish that mobility and atomicity
matter. Our focus is the interaction among routing, the measured cost of
moving a live trajectory, replay-equivalent state validation, and
exactly-once external effects. The effect protocol is inspired by classic
compensating transactions \citep{garciamolina1987sagas}, but it is not a
distributed database transaction: external tools may offer only idempotency
keys or compensating actions, and \system{} records that limitation.

\paragraph{Agent benchmarks.}
ALFWorld evaluates text agents grounded in interactive household tasks
\citep{shridhar2021alfworld}; WebShop evaluates compositional web interaction
\citep{yao2022webshop}; and AppWorld provides executable APIs and stateful
application tasks \citep{trivedi2024appworld}. AgentProcessBench evaluates
intermediate process quality rather than only final answers
\citep{fan2026agentprocessbench}. We use these public resources for distinct
roles: ALFWorld and WebShop measure long-horizon routing; AppWorld stresses
effectful calls; AgentProcessBench provides a process-level diagnostic.

\section{\system{} Design}
\label{sec:design}

\subsection{Design goals and non-goals}

The protocol has four goals. First, a receiver either reconstructs a state
that satisfies declared invariants or refuses execution. Second, transfer
costs remain observable at component level. Third, a repeated handoff or tool
retry cannot silently duplicate a committed effect. Fourth, all learned
estimators are train-split artifacts and can be replaced without weakening
protocol checks.

\system{} does not guarantee semantic equivalence between arbitrary natural
language summaries, global serializability across unrelated tools, or success
from an incapable base model. It also does not infer a hidden world state that
the benchmark does not expose. These boundaries are important: a protocol can
prevent structurally invalid continuation and duplicated effects without
proving that an LLM's plan is correct.

\subsection{Relay state packet}
\label{sec:packet}

Table~\ref{tab:packet} summarizes the packet. Canonical JSON serialization
sorts keys and uses normalized separators before hashing. Each packet contains
a schema version, immutable task identity, sequence number, parent-state hash,
payload hash, provenance records, declared invariants, and effect records.
The sealed packet hash covers every field required for validation.

\begin{table}[t]
\centering
\small
\caption{Logical contents of a relay state packet. All identifiers and hashes
are deterministic under canonical serialization.}
\label{tab:packet}
\begin{tabular}{p{0.20\linewidth}p{0.72\linewidth}}
\toprule
Field & Purpose \\
\midrule
Task identity & Benchmark, split, immutable example identifier, and task hash \\
Version chain & Schema version, sequence number, parent packet hash, sender, receiver \\
World state & Normalized observation and environment-visible facts \\
Evidence & Tool outputs or observations with source, timestamp, and content hash \\
Plan state & Current objective, completed steps, open obligations, and constraints \\
Effect ledger & Canonical effect key, class, status, attempt, result hash, compensation link \\
Integrity & Required paths, payload hash, packet hash, and provenance declarations \\
\bottomrule
\end{tabular}
\end{table}

The packet separates \emph{observed facts} from model-generated plan state.
Evidence items retain their origin and hash; generated rationales cannot be
silently relabeled as environment observations. Task identity binds a packet
to one immutable benchmark instance, preventing accidental reuse across
episodes. A monotonically increasing sequence and parent hash make stale or
forked updates detectable.

\subsection{Transfer modes and typed deltas}
\label{sec:delta}

The compound action chooses one of three transfer modes:
\begin{description}
  \item[\methodtag{Reuse}] The target executor already holds the exact packet
  hash, so no payload is transmitted.
  \item[\methodtag{Full}] The sender transmits the complete sealed packet.
  This is the fidelity reference and the fallback when no common parent exists.
  \item[\methodtag{Delta}] The sender transmits typed patch operations against
  a receiver-known parent hash.
\end{description}

A delta $d_t=(h_{t-1},\Delta_t,h_t)$ names its parent and expected child
hashes. Operations are restricted to typed paths and explicit
\textsc{add}, \textsc{replace}, or \textsc{remove} semantics. Applying a
delta is accepted only when the local parent hash equals $h_{t-1}$ and the
materialized child seals to $h_t$. This design prevents a common ambiguity in
free-form summaries: absence can mean either unchanged, deliberately removed,
or accidentally omitted. The operation type makes those cases distinct.

Content-addressed trace storage de-duplicates packets and evidence blobs.
Delta selection compares the estimated serialized size and validation cost
with full transfer; it does not assume that a delta is always smaller. If the
parent is absent or the estimated delta is larger, \methodtag{Full} remains a
valid candidate.

\subsection{Receiver validation and patch-on-demand}
\label{sec:validation}

Before invoking a model, the receiver runs deterministic checks in the
following order:
\begin{enumerate}
  \item verify schema version, canonical payload hash, and packet hash;
  \item match benchmark, split, example identifier, and task hash;
  \item verify the declared parent and monotone sequence;
  \item check required state paths and evidence provenance;
  \item reject duplicate effect keys or illegal effect-state transitions; and
  \item evaluate task-adapter invariants that depend only on public
  environment state.
\end{enumerate}

Failure is fail-closed. If the missing information is patchable, the receiver
returns a structured request naming required paths and the expected parent.
The sender responds with a sealed patch bundle, after which the receiver
re-validates the full materialized state. Hash mismatch, task mismatch, or an
illegal effect transition is not patchable and aborts the handoff. Patch
frequency and patch bytes are included in the relay tax rather than removed
from latency reports.

We define \emph{state fidelity} relative to full replay. Given the same model
commit, decoding configuration, random seed, environment state, and step,
a transferred state is replay-equivalent when it passes all deterministic
invariants and produces the same canonical next tool action as full replay.
This is an operational measure, not a claim of equivalence between hidden
model representations.

\subsection{Effect ledger and commit modes}
\label{sec:ledger}

Each externally visible tool action is normalized into an intent tuple
$(\mathrm{task},\mathrm{tool},\mathrm{arguments},\mathrm{scope})$ and assigned
a canonical effect key. Tool adapters classify actions as read-only,
idempotent, compensatable, or irreversible. A ledger record moves through
\textsc{authorized}, \textsc{prepared}, \textsc{committed}, and, where
supported, \textsc{compensated}. The transition relation rejects a second
commit for the same key and rejects a different normalized intent that
collides with an existing key.

The policy chooses a commit mode compatible with the adapter:
\begin{description}
  \item[\methodtag{Read}] executes side-effect-free calls without a commit;
  \item[\methodtag{Idempotent}] passes the canonical key to a tool that
  natively de-duplicates retries;
  \item[\methodtag{Prepare--commit}] records authorization and preparation
  before the external call, then commits the returned result hash; and
  \item[\methodtag{Compensating}] additionally records the inverse action
  when atomic rollback is unavailable.
\end{description}

An uncertain timeout never becomes an implicit retry. The adapter first
queries by idempotency key when possible; otherwise the episode is marked
indeterminate and handled according to the benchmark's official semantics.
This conservative rule may reduce nominal completion rate, but avoids
labeling duplicated real effects as success.

\subsection{Joint action selection}
\label{sec:policy}

For every valid $a=(e,m,c)$, the selector estimates success probability,
inference cost, relay-tax components, latency, and protocol risk. Candidate
features include task metadata, trajectory length, tool class, recent
failures, packet sizes, receiver cache state, and the current recorded link
sample. The feature schema is fixed before evaluation. Separate probability
and cost heads are fit only on native rollouts from the official training
split; validation selects regularization and the quality threshold.

Selection is constrained rather than purely score-based. Invalid commit modes
and memory-infeasible executors are removed. Among actions meeting the
validation-chosen success and latency constraints, the policy minimizes
predicted measured cost. If none is feasible, a declared fallback chooses the
highest predicted success among protocol-valid actions. Ties are deterministic.
For analysis, the same predictions are passed to a zero-tax counterfactual
that sets $H_t(a)=0$. Its disagreements with the deployed selector identify
candidate switch-cost reversals without changing executed trajectories.

This formulation couples routing and representation. A large model with a
full replay, the same model with a delta, and staying local with resident
state are different candidates. A sequential baseline that first chooses a
model and then applies one fixed transfer rule cannot express this trade-off.

\section{Implementation}
\label{sec:implementation}

The research artifact is a Python package with no mandatory machine-learning
dependency in its protocol core. The core defines packet schemas, canonical
hashing, state diff and application, validation, the effect ledger, measured
cost accounting, constrained selection, provenance manifests, metrics, and
statistical utilities. This separation allows protocol and failure tests to
run without downloading a model.

Native inference is implemented through pinned Hugging Face Transformers
revisions using each model's official chat template. The executor records
input and output tokens, wall-clock time, peak device memory, prompt hash,
model commit, tokenizer commit, decoding parameters, and seed. Model outputs
are immutable run artifacts. Benchmark adapters, rather than prompts, expose
official observations, actions, termination conditions, and evaluation
results.

The initial joint estimator is intentionally lightweight. Per-action logistic
heads estimate success and ridge heads estimate continuous costs from a fixed
feature schema. Training rows are generated by native rollouts on official
training tasks. The package rejects rows whose provenance indicates a
validation or test split. Learned artifacts store their training-manifest
hash so that a result can be traced back to data, code, model, and command.
This estimator is not presented as a learning contribution; it is a controlled
way to test the systems mechanisms.

Configuration is fail-closed. A formal configuration requires immutable model,
dataset, and repository commits, disallows a sample limit, enables native
inference, disables test-label access, and fixes all storage beneath
\texttt{/root/autodl-tmp/AgentRelay}. A local configuration can run bounded
diagnostics but is explicitly marked as non-evidence. Public benchmark
repositories and Hugging Face datasets are downloaded only after revisions
are resolved to immutable commits.

\section{Evaluation Methodology}
\label{sec:evaluation}

\subsection{Research questions}

The evaluation is organized around four primary questions:
\begin{description}
  \item[RQ1:] How large is the relay tax, and how often does measuring it
  reverse an executor or payload decision?
  \item[RQ2:] Does joint executor--payload--commit selection improve the
  Pareto frontier relative to fixed transfer or model-only routing?
  \item[RQ3:] Do typed deltas with validation and patching reduce transfer
  while retaining replay-equivalent continuation state?
  \item[RQ4:] Does the effect ledger prevent duplicated or conflicting
  mutations under handoff failures and retries?
\end{description}
Each question has a separate experiment and rejection condition; aggregate
task success alone cannot validate all four.

\subsection{Public benchmarks and immutable splits}
\label{sec:benchmarks}

All task instances come from public benchmark releases. We neither generate
new tasks nor paraphrase existing ones.

\paragraph{ALFWorld.}
ALFWorld contains 3,827 text-interactive household tasks
\citep{shridhar2021alfworld}. Training rollouts use only its official training
split. Formal evaluation reports the complete standard \texttt{valid\_seen}
and \texttt{valid\_unseen} partitions separately and jointly (140 and 134
tasks in the pinned distribution). Episodes are capped at 50 environment
steps, with 2,048 prompt tokens and 512 response tokens.

\paragraph{WebShop.}
WebShop contains 12,087 crowd-sourced instructions grounded in a corpus of
approximately 1.18 million products \citep{yao2022webshop}. We retain the
official train, development, and test manifests from the pinned release.
Formal evaluation uses the complete official test manifest, a 20-step cap,
4,096 prompt tokens, and 1,024 response tokens.

\paragraph{AppWorld.}
AppWorld provides stateful application APIs and executable evaluators
\citep{trivedi2024appworld}. We retain its official 105/60/168/417
train/dev/test-normal/test-challenge split. Training and trace generation use
the 105 training tasks; development selects thresholds; both test partitions
are evaluated once by the frozen pipeline. Episodes use a 50-step cap, 4,096
prompt tokens, and 1,024 response tokens. Mutation-safety claims are confined
to this official sandbox.

\paragraph{Diagnostic data.}
AgentProcessBench provides 1,000 public trajectories spanning BFCL, GAIA-dev,
HotpotQA, and tau2 \citep{fan2026agentprocessbench}. It is used only for local
schema and process diagnostics. Its labels neither train the final main-task
router nor score the main comparison.

\paragraph{Network conditions.}
Communication robustness replays payloads through ordered, public recorded
cellular traces, beginning with the Verizon LTE traces distributed with
Mahimahi \citep{netravali2015mahimahi}. Samples retain their original order;
traces are never randomized, interpolated into synthetic tasks, or repeated
past exhaustion. Trace-driven communication latency is reported separately
from live GPU wall-clock latency.

\subsection{Models and hardware}

The primary pair uses Gemma 4 instruction-tuned models
\citep{gemma2026gemma4}: \texttt{google/gemma-4-E4B-it} as the edge executor
and \texttt{google/gemma-4-12b-it} as the cloud-equivalent executor, each at an
immutable revision and with the declared memory-safe loading precision. A shared family reduces tokenizer and prompt-format
confounds while retaining a substantial capacity difference. Revisions are
resolved and locked before any evidence run. Greedy decoding is the default;
if an official harness requires sampling, three declared seeds are used
identically across paired methods.

The local RTX 4080 Laptop is restricted to unit tests, dependency validation,
short public-data pipeline checks, and a small non-reportable pilot. All
router fitting, complete benchmark inference, comparisons, ablations, and
reported measurements run on an AutoDL RTX 4090D. The two models are profiled
as warm services sequentially when simultaneous residency would introduce a
memory confound. Replayed service time is labeled as such and never presented
as simultaneous live execution.

\subsection{Baselines and fairness controls}
\label{sec:baselines}

We compare edge-only, cloud-only, random routers, a FrugalGPT-style cascade
\citep{chen2024frugalgpt}, a RouteLLM-style task router
\citep{ong2025routellm}, an executor-only step router, an AgentRouter-style
classifier \citep{paul2026agentrouter}, and a Hera-style device--cloud router
\citep{zhang2026hera}. Official implementations are used when they support
the pinned benchmark and model pair. Otherwise we label the method as a
paper-faithful reimplementation and publish every deviation; we do not paste
numbers from incompatible papers into the main table.

State-transfer baselines are full replay, recency truncation, a frozen
narrative-summary prompt, an unverified structured snapshot, typed delta
without repair, and the full verify-and-patch protocol. Narrative and
structured baselines are matched to the typed payload by target-token count
where feasible, and unmatched cases are disclosed. Reliability baselines are
no ledger, key-based de-duplication only, and an Atomix-inspired progress
barrier \citep{mohammadi2026atomix}. All methods share task prompts, tool
schemas, model commits, decoding, maximum tokens, environment-step caps, task
order, seeds, and hardware profiles.

An oracle that sees realized per-action outcomes may be reported only as a
non-deployable upper bound. It is excluded from statistical comparisons and
never supplies labels, features, or thresholds to a deployed method.

\subsection{Metrics and statistical protocol}
\label{sec:metrics}

Task quality uses each benchmark's official success or reward, supplemented
by invalid-action rate and environment steps. Efficiency reports end-to-end
trajectory latency, median and 95th-percentile step latency, cloud-step ratio,
switches per trajectory, transfer bytes, target-tokenized input, peak GPU
memory, and each component of Equation~\ref{eq:relay-tax}. Warm and cold model
loading are separate conditions.

State fidelity reports required-invariant recall, hash and provenance
validity, stale-version rejection, first-pass rehydration, patch rate and
rounds, patch bytes, canonical next-action agreement with full replay, and
continuation success. Effect correctness reports duplicate commits,
conflicting mutations, official final-state correctness, collateral state
failures, recovery success, and recovery latency.

Methods see tasks in identical order with identical seeds. Binary outcomes are
paired by task and analyzed with exact McNemar tests; continuous paired
differences use deterministic task-level bootstrap confidence intervals.
Primary comparisons apply Holm correction. For stochastic components we
report three independent runs, mean and standard deviation for convention,
and task-level 95\% paired intervals. Heavy-tailed latency and size are
summarized by median and p95 rather than a normal approximation. Exact
commands, raw JSONL records, manifests, and correction families are retained.

\subsection{Experiment matrix and predeclared decisions}

\begin{table}[t]
\centering
\small
\caption{Predeclared experiments. A claim is supported only if its primary
test and mechanism-specific diagnostics agree.}
\label{tab:experiments}
\begin{tabular}{p{0.08\linewidth}p{0.39\linewidth}p{0.20\linewidth}p{0.22\linewidth}}
\toprule
ID & Intervention & Primary outcome & Claim decision \\
\midrule
E1 & Full frozen-policy comparison & Success and Pareto frontier & C2 \\
E2 & Forced matched-position handoffs & Fidelity and transfer size & C3 \\
E3 & Zero-tax versus measured-tax replay & Validated route reversals & C1 \\
E4 & Four interruption points in AppWorld & Duplicate effects and final state & C4 \\
E5 & Ordered recorded network traces & Frontier by trace condition & C1--C2 robustness \\
E6 & Same-pair cross-benchmark transfer & Directional replication & Scope only \\
E7 & Packet and controller scaling & Latency, bytes, memory & Systems overhead \\
\bottomrule
\end{tabular}
\end{table}

E2 forces handoffs at positions selected from the manifest without outcome
labels. E3 replays the same frozen states through a zero-tax counterfactual,
an inference-only cost model, and the measured full objective, then validates
changed choices by native continuation. E4 interrupts immediately before
tool send, after send before acknowledgement, after acknowledgement before
packet commit, and during target retry. Failure injection changes execution
timing only; it does not alter tasks, desired answers, or evaluator state.

The claim criteria are conservative. C1 requires nonzero validated routing
reversals and an improved realized objective, not merely measurable overhead.
C2 requires nondomination or a statistically supported improvement in a
predeclared quality-constrained region, not one favorable scalarization. C3
requires lower transfer than full replay while preserving both invariant and
continuation fidelity; smaller packets alone are insufficient. C4 requires
zero duplicate committed effects under the declared injections and no
degradation of official final-state correctness relative to the strongest
runnable reliability baseline.

\section{Results}
\label{sec:results}

No formal result exists at manuscript-drafting time. Tables
\ref{tab:main-results}--\ref{tab:ablation-results} are schemas connected to
the immutable result pipeline; \tbd{} is not a zero and must not be included
in an average, plot, or claim.

\begin{table*}[t]
\centering
\small
\caption{Main frozen-policy comparison on complete official evaluation
splits. Values will include 95\% paired intervals where applicable.
APB is diagnostic-only and is excluded from this table.}
\label{tab:main-results}
\resizebox{\textwidth}{!}{
\begin{tabular}{llrrrrrr}
\toprule
Benchmark & Method & Success or reward & E2E latency & Cloud steps & Switches &
Transfer bytes & Peak GPU memory \\
\midrule
ALFWorld seen/unseen & Edge-only & \tbd & \tbd & \tbd & \tbd & \tbd & \tbd \\
 & Cloud-only & \tbd & \tbd & \tbd & \tbd & \tbd & \tbd \\
 & Model-only step router & \tbd & \tbd & \tbd & \tbd & \tbd & \tbd \\
 & Hera-style router & \tbd & \tbd & \tbd & \tbd & \tbd & \tbd \\
 & \system{} & \tbd & \tbd & \tbd & \tbd & \tbd & \tbd \\
\midrule
WebShop test & Edge-only & \tbd & \tbd & \tbd & \tbd & \tbd & \tbd \\
 & Cloud-only & \tbd & \tbd & \tbd & \tbd & \tbd & \tbd \\
 & Model-only step router & \tbd & \tbd & \tbd & \tbd & \tbd & \tbd \\
 & Hera-style router & \tbd & \tbd & \tbd & \tbd & \tbd & \tbd \\
 & \system{} & \tbd & \tbd & \tbd & \tbd & \tbd & \tbd \\
\midrule
AppWorld normal/challenge & Edge-only & \tbd & \tbd & \tbd & \tbd & \tbd & \tbd \\
 & Cloud-only & \tbd & \tbd & \tbd & \tbd & \tbd & \tbd \\
 & Model-only step router & \tbd & \tbd & \tbd & \tbd & \tbd & \tbd & \tbd \\
 & Hera-style router & \tbd & \tbd & \tbd & \tbd & \tbd & \tbd \\
 & \system{} & \tbd & \tbd & \tbd & \tbd & \tbd & \tbd \\
\bottomrule
\end{tabular}}
\end{table*}

\begin{table}[t]
\centering
\small
\caption{Relay-tax decomposition and routing reversal analysis (E3).}
\label{tab:relay-results}
\begin{tabular}{lrrr}
\toprule
Quantity & ALFWorld & WebShop & AppWorld \\
\midrule
Encode / E2E latency & \tbd & \tbd & \tbd \\
Transfer / E2E latency & \tbd & \tbd & \tbd \\
Verify and patch / E2E & \tbd & \tbd & \tbd \\
Rehydration / E2E & \tbd & \tbd & \tbd \\
Decision reversal rate & \tbd & \tbd & \tbd \\
Validated beneficial reversals & \tbd & \tbd & \tbd \\
\bottomrule
\end{tabular}
\end{table}

\begin{table}[t]
\centering
\small
\caption{Matched-position state-transfer comparison (E2).}
\label{tab:fidelity-results}
\resizebox{\linewidth}{!}{
\begin{tabular}{lrrrr}
\toprule
Transfer mode & Bytes & Invariant recall & Action agreement & Continuation success \\
\midrule
Full replay & \tbd & \tbd & \tbd & \tbd \\
Recency truncation & \tbd & \tbd & \tbd & \tbd \\
Narrative summary & \tbd & \tbd & \tbd & \tbd \\
Unverified snapshot & \tbd & \tbd & \tbd & \tbd \\
Typed delta, no patch & \tbd & \tbd & \tbd & \tbd \\
\system{} verify and patch & \tbd & \tbd & \tbd & \tbd \\
\bottomrule
\end{tabular}}
\end{table}

\begin{table}[t]
\centering
\small
\caption{AppWorld interruption experiment (E4). Counts retain their
task-level denominators in the result artifact.}
\label{tab:effect-results}
\begin{tabular}{lrrr}
\toprule
Method & Duplicate commits & Correct final state & Recovery latency \\
\midrule
No ledger & \tbd & \tbd & \tbd \\
De-duplication only & \tbd & \tbd & \tbd \\
Commit barrier & \tbd & \tbd & \tbd \\
\system{} & \tbd & \tbd & \tbd \\
\bottomrule
\end{tabular}
\end{table}

\begin{table}[t]
\centering
\small
\caption{Mechanism ablations aggregated only after reporting
benchmark-specific results.}
\label{tab:ablation-results}
\resizebox{\linewidth}{!}{
\begin{tabular}{lrrrr}
\toprule
Variant & Success & Measured cost & Transfer bytes & Invalid or duplicate effects \\
\midrule
Full \system{} & \tbd & \tbd & \tbd & \tbd \\
Without relay tax & \tbd & \tbd & \tbd & \tbd \\
Executor-only action & \tbd & \tbd & \tbd & \tbd \\
Fixed typed delta & \tbd & \tbd & \tbd & \tbd \\
Without validation & \tbd & \tbd & \tbd & \tbd \\
Verify without patch & \tbd & \tbd & \tbd & \tbd \\
Without effect ledger & \tbd & \tbd & \tbd & \tbd \\
\bottomrule
\end{tabular}}
\end{table}

\paragraph{Planned interpretation.}
After execution, this paragraph will report effect sizes and uncertainty
before significance labels. A route reversal counts as beneficial only when
the native continuation realizes a lower constrained objective. Pareto claims
will display all measured operating points; a single hand-selected threshold
will not stand in for a frontier. State failures will be categorized as
missing field, stale parent, hash or provenance mismatch, patch exhaustion,
or continuation disagreement. Effect failures will identify the interruption
point and tool class. Negative and null outcomes remain in the paper.

\section{Discussion}
\label{sec:discussion}

\subsection{What each possible outcome would mean}

The design admits informative negative results. If the measured relay tax is
small and E3 finds no validated reversal, C1 is rejected; the paper should
then present the protocol as a state-correctness mechanism rather than a
routing advance. If joint selection matches an executor-only policy, C2 is
rejected even if deltas save bytes. If typed deltas are compact but reduce
continuation fidelity, C3 is rejected and full replay becomes the appropriate
default. If the ledger fails at an interruption point, C4 is restricted to
the tool classes and transitions that pass. These outcomes are encoded in the
project gates so that wording follows evidence instead of retroactively
changing the experiment.

Several combinations are also diagnostically useful. A high relay tax with no
beneficial reversal would imply estimator or constraint error rather than an
irrelevant systems cost. Strong invariant recall with weak next-action
agreement would indicate that the declared schema omits decision-relevant
state. Zero duplicate commits with poor task success could mean that the
ledger is safe but too conservative. Reporting these mechanism-level signals
prevents one end-to-end number from masking the cause.

\subsection{Deployment implications}

The action space makes caching and network conditions visible to policy.
Resident state can favor staying on an otherwise weaker executor; a small
delta can favor switching; and an irreversible effect can force a commit
barrier even when it increases latency. This is relevant whenever model
placement and state placement differ. It also suggests a practical interface
between agent frameworks and routers: the framework should expose packet
availability, effect class, and measurable handoff components instead of
offering only prompt text and model scores.

The protocol deliberately defaults unknown effect classes to the safer
barrier path. This is appropriate for benchmark and enterprise automation but
can be too conservative for high-throughput, low-impact tools. Deployments
may define adapter-specific risk policies, provided that those policies are
fixed before evaluation and that unknown effects are not silently treated as
read-only.

\section{Limitations and Threats to Validity}
\label{sec:limitations}

\paragraph{Benchmark scope.}
Three interactive benchmark families cannot represent every long-running
agent. ALFWorld and WebShop emphasize textual interaction, whereas AppWorld
offers a sandbox rather than irreversible real services. AgentProcessBench is
diagnostic only. Conclusions must therefore be scoped to the tested adapters,
task distributions, and public releases.

\paragraph{Exactly-once boundary.}
Exactly-once behavior is attainable only when the tool accepts an idempotency
key, exposes queryable effect state, or supports a reliable compensation.
With a non-queryable irreversible service and an acknowledgement lost after
execution, no client-side protocol can distinguish an unexecuted action from
an executed but unacknowledged action. \system{} labels that case
indeterminate and refuses an automatic retry. The paper must not generalize
sandbox results to uncooperative external APIs.

\paragraph{State fidelity.}
Canonical action agreement is an observable proxy for continuation fidelity,
not proof that two model executions have identical internal states. The typed
schema may omit a fact that becomes relevant later, and benchmark-specific
invariants may be incomplete. Full replay is therefore retained as a fallback
and reference rather than declared unnecessary.

\paragraph{Hardware realism.}
A single 4090D can reproduce native inference and warm service profiles but is
not a physical edge--cloud deployment. Recorded public traces isolate
communication sensitivity, while sequential model profiling avoids
contention; neither captures every production queue, network stack, or remote
accelerator. We report live, profiled, and trace-replayed quantities
separately. Claims about absolute production latency require external
validation.

\paragraph{Model and estimator dependence.}
The primary same-family model pair controls some prompt-format variation but
may favor transfer compatibility. The lightweight estimator can be
miscalibrated under distribution shift. A same-pair cross-benchmark replication
tests directionality across workloads, yet does not establish universality. Test labels remain
unavailable to the router, so calibration failures cannot be repaired after
test evaluation.

\paragraph{Instrumentation effects.}
Canonical serialization, hashing, and validation add controller work. This
cost is part of the treatment and is measured, but fine-grained timing itself
can perturb short operations. E7 repeats microbenchmarks, separates warm-up,
and retains raw samples. Result interpretation emphasizes end-to-end
differences and intervals rather than sub-millisecond point estimates.

\paragraph{Security and privacy.}
Integrity hashes detect accidental or adversarial modification but do not
encrypt state or authenticate principals. Evidence packets may contain
sensitive observations in a real deployment. Encryption, key management,
access control, retention, and traffic-analysis resistance are outside this
paper and are required before handling private data.

\section{Responsible Research and Reproducibility}
\label{sec:responsible}

The study uses public benchmark tasks and sandboxed environments; it does not
collect human-subject data or invoke real purchases, messages, or account
mutations. Repository and dataset licenses are recorded at download time.
Benchmark text is not modified, and no synthetic task augments an evaluation
split. Narrative examples, if later included, are selected by a declared
representative, difficult, or failure rule and retain their public sample
identifiers.

Every formal run emits a manifest containing dataset and model commits,
sample identifiers, split, seed, prompt and code hashes, configuration,
command, hardware, timestamp, and output hashes. Raw native outputs are
append-only. Aggregation consumes those records and does not accept manually
entered result cells. The locked formal configuration requires all models,
datasets, repositories, logs, and outputs below
\texttt{/root/autodl-tmp/AgentRelay}, which persists across AutoDL instance
restarts. GPU hours, peak memory, failed runs, and exclusions will accompany
the final artifact.

\section{Conclusion}
\label{sec:conclusion}

Edge--cloud agent routing is incomplete when it treats a switch as only a
model-selection decision. \system{} represents the switch as a measured,
validated, and effect-aware transaction. Its four proposed innovations are a
relay-tax-aware objective, joint executor--payload--commit selection, typed
deltas with verify-and-patch recovery, and an idempotent effect ledger. The
artifact and evaluation plan isolate each mechanism using public tasks,
immutable splits, native inference, recorded network traces, and paired
statistics. Whether these mechanisms improve real quality--cost--latency and
correctness trade-offs remains an empirical question; this draft reserves
that conclusion until the formal 4090D runs populate the registered result
artifacts.

\appendix

\section{Claim--Evidence Contract}
\label{app:claims}

\begin{table*}[h]
\centering
\small
\caption{Contract between manuscript claims and required evidence. No claim
graduates from proposed to supported without all required rows.}
\label{tab:claim-contract}
\begin{tabular}{p{0.06\textwidth}p{0.26\textwidth}p{0.30\textwidth}p{0.28\textwidth}}
\toprule
Claim & Mechanism & Required evidence & Failure or narrowing rule \\
\midrule
C1 & Component-measured relay tax in routing &
E3 component timings, zero-tax disagreements, and native validation of changed
choices & Reject if no beneficial validated reversal; retain overhead profile \\
C2 & Joint executor, transfer, and commit action &
E1 complete-split frontier plus executor-only and fixed-transfer comparisons &
Reject if frontier is dominated or indistinguishable in the declared region \\
C3 & Typed delta, deterministic validation, narrow patch &
E2 matched positions and size, invariant, action, and continuation metrics &
Reject if transfer is not lower or fidelity falls beyond the validation bound \\
C4 & Canonical effect keys and legal state transitions &
E4 all four interruption points, official final state, duplicates, and recovery &
Restrict to passing tool classes; reject broad exactly-once wording on any
unexplained duplicate \\
\bottomrule
\end{tabular}
\end{table*}

\section{Handoff Procedure}
\label{app:procedure}

\paragraph{Sender.}
The sender seals the current packet, checks whether the receiver acknowledges
its hash, enumerates protocol-valid compound actions, and records the selected
action and all predictions. For a delta it materializes a typed diff against
the acknowledged parent. For full transfer it sends the complete canonical
packet. The sender never removes a committed effect record.

\paragraph{Receiver.}
The receiver obtains the named parent from content-addressed storage, applies
the delta if present, recomputes hashes, and runs the validation sequence in
Section~\ref{sec:validation}. A patchable failure returns exact paths or trace
references. Other failures abort. Only a validated packet is rendered through
the executor's official chat template.

\paragraph{Tool boundary.}
The adapter canonicalizes an intended call before execution. Read-only calls
are logged. A mutating call must be authorized under a compatible commit mode,
and its key is transferred with state. On retry, a committed key returns the
recorded result when available; a conflicting intent is rejected. An
uncertain unqueryable effect becomes indeterminate rather than automatically
repeated.

\section{Formal Execution Checklist}
\label{app:checklist}

\begin{enumerate}
  \item Resolve every requested Hugging Face and Git revision to an immutable
  commit; record upstream hashes for any dataset mirror.
  \item Run protocol unit tests and configuration validation. Confirm the
  formal root, no sample limit, native inference, and disabled test access.
  \item Download official assets to the persistent AutoDL data disk. Retain
  licenses and unmodified manifests.
  \item Generate train-split native rollouts, fit estimator heads, and select
  thresholds on development data only.
  \item Freeze model, prompt, adapter, policy, threshold, trace, and seed
  manifests. Hash the complete formal configuration.
  \item Execute E1--E4 on complete official partitions in fixed task order.
  Append E5--E7 only under their declared dependencies and budget.
  \item Aggregate JSONL records mechanically, run paired uncertainty tests and
  Holm correction, and populate tables from the emitted machine-readable
  summary.
  \item Audit every numeric statement against its result record and every
  citation against its bibliographic source before removing the
  pre-experimental banner.
\end{enumerate}

\section{Reporting Fields}
\label{app:reporting}

For each method--benchmark cell, the result artifact contains the denominator,
task identifiers, seeds, success or reward, environment steps, invalid calls,
edge and cloud tokens, per-component latency, packet and patch bytes, switches,
packet validation status, effect transitions, peak memory, and exit reason.
Exclusions require a machine-readable reason and remain visible in the
denominator report. Crash and timeout rates are reported rather than silently
retried until success.

For each learned estimator, the artifact additionally records the fixed
feature schema, training split and trace-manifest hashes, regularization,
calibration error on development data, chosen constraints, and all candidate
predictions at decision time. This permits off-policy inspection of routing
reversals without granting the policy access to realized test outcomes.

\bibliographystyle{plainnat}
\bibliography{references}

\end{document}
